\begin{figure}[h!]
    \tdplotsetmaincoords{65}{20}
    \centering
    \begin{tikzpicture}[tdplot_main_coords, scale=1.0]

        % Dimensions
        \draw[thick,->] (0, 0, 0) -- (4.5, 0, 0) node[anchor=west]{$\log(r)$};
        \draw[thick,->] (0, 0, 0) -- (0, 0, 2) node[anchor=south]{$\log(s)$};
        \draw[thick,->] (4, 0, 0) arc (0:90:4) node[midway, anchor=south west]{$\phi$};

        % Adaptive phi grid
        \begin{scope}[canvas is xy plane at z=0, opacity=0.2]
            % Arcs
            \draw (1,0,0) arc (0:90:1);
            \draw (2,0,0) arc (0:90:2);
            \draw (3,0,0) arc (0:90:3);
            % \draw (4,0,0) arc (0:90:4);

            % Level 1
            \draw (45:1) -- (45:4);

            % Level 2
            \draw (22.5:2) -- (22.5:4);
            \draw (67.5:2) -- (67.5:4);

            % Level 3
            \draw (11.25:3) -- (11.25:4);
            \draw (33.75:3) -- (33.75:4);
            \draw (56.25:3) -- (56.25:4);
            \draw (78.75:3) -- (78.75:4);

            % Axes
            \draw[->] (0, 0, 0) -- (0, 4.5, 0);
            % \draw[->] (0, 0, 0) -- (4.5, 0, 0);
        \end{scope}

        % Heptahedron
        \begin{scope}[thick, color=violet]
            % Bottom
            \begin{scope}[canvas is xy plane at z=0]
                \draw (45:2) -- (45:3) arc [start angle=45, end angle=67.5, radius=3] -- (67.5:2) arc [start angle=67.5, end angle=45, radius=2];
                \coordinate (A) at (45:2);
                \coordinate (B) at (45:3);
                \coordinate (C) at (56.25:3);
                \coordinate (D) at (67.5:3);
                \coordinate (E) at (67.5:2);
                \coordinate (P0) at (51:2.6);
            \end{scope}

            \begin{scope}[canvas is xy plane at z=0.2]
                \coordinate (P) at (51:2.6);
            \end{scope}

            % Top
            \begin{scope}[canvas is xy plane at z=0.5]
                \draw (45:2) -- (45:3) arc [start angle=45, end angle=67.5, radius=3] -- (67.5:2) arc [start angle=67.5, end angle=45, radius=2];
                \coordinate (F) at (45:2);
                \coordinate (G) at (45:3);
                \coordinate (H) at (56.25:3);
                \coordinate (I) at (67.5:3);
                \coordinate (J) at (67.5:2);
                \coordinate (P1) at (51:2.6);
            \end{scope}

            \draw (A) -- (F);
            \draw (B) -- (G);
            \draw (C) -- (H);
            \draw (D) -- (I);
            \draw (E) -- (J);
            \draw (P0) circle (0.02);
            \draw (P) circle (0.02) node[left, inner sep=1pt]{\tiny$\vec{p}$};
            \draw (P1) circle (0.02);
            \draw[thin, densely dotted] (P0) -- (P1);
        \end{scope}
    
    \end{tikzpicture}
    \vspace{-5pt}
    \caption{
        The footprint shape $\vec{p}$ is discretized using logarithmic steps along the $r$ and $s$ dimension and adaptive step sizes along the angular $\phi$ dimension. As a result, the grid cells take on a heptahedral shape (violet).
        \label{fig:footprint-grid}
    }
\end{figure}